\abstract{
Graph-based traffic forecasting methods typically rely on predefined adjacency matrices derived from physical road proximity. However, dense physical connectivity can induce excessive spatial aggregation---a form of oversmoothing---that degrades prediction accuracy. We propose a multi-lag graph structure learning framework that discovers temporally heterogeneous functional dependencies between traffic sensors. Using DAGMA continuous optimization, we construct lag-specific dependency graphs from explicit temporal input blockings of the form $Z = [x(t{-}L), \ldots, x(t)]$, yielding separate functional graphs for each temporal lag. A T-GCN-MultiGSL-Mix architecture then combines these lag-specific graphs through a per-node, per-timestep learned mixing gate. Experiments on the Los-loop highway dataset demonstrate a mean RMSE improvement of 14.9\% over a no-graph baseline across five random seeds at 5-minute sampling, increasing to 27.4\% at 15-minute sampling, with the improvement robust to a parameter-matched control and to matched-edge-budget sparse-graph controls. The original submission's central result---a sparse learned DAG substantially outperforming the dense physical graph---is reproduced under the revised protocol, with a mean improvement of 21.7\% across five seeds and the historically submitted graphs recovered at the support level. A lag ablation study confirms that all three temporal lags contribute complementary information. On the SZ-Taxi urban dataset, improvements are marginal, indicating dataset dependence. These findings establish that temporally heterogeneous graph structures can substantially benefit traffic forecasting when the underlying dependency patterns vary across temporal scales.
}
